% !TEX root = ../linnik399.tex
\section{The graded near-density lemma}\label{sec:near}

The purpose of this section is to constrain many characters at once while retaining
the location of each selected zero. We first bound a quadratic form in smoothed prime
sums, called \emph{responses} (\cref{thm:graded}). We then bound each response from below
using selected zeros (\cref{lem:response}). Finally, an elementary optimization argument
introduces the single threshold used in a near row (\cref{prop:threshold}).
The analytic inputs are stated in \cref{sec:inputs}; the scope of the formalization is described
in \cref{subsec:lean}. \Cref{subsec:nearidea} explains the idea of the proof, and
\cref{ex:twolevels} shows what the resulting constraint can do that separate counts cannot.

\subsection{Data}\label{subsec:neardata}
The detector determines the response to a zero. The Gram test and sieve weight control
correlations between characters. The anchors may vary between entries, but all test
functions and allowed offsets are fixed before $q$ tends to infinity.

\begin{definition}[Near-density data]\label{def:neardatum}
A \emph{near datum} consists of the following fixed objects; Conditions~1 and~2 are
those of \cref{def:conditions}.
\begin{itemize}
  \item A \emph{detector} $f$ satisfying Condition~1, with transform $F$.
  \item A \emph{Gram test} $g$ satisfying Conditions~1 and~2, with transform $G$.
  \item Real numbers $s_1$ and $s$, called the \emph{safe anchor} and the
    \emph{reference anchor}.
  \item A \emph{sieve weight}: cells $t_0<t_1<\dots<t_m$ with $t_0>\frac13+2\eps'$ for a fixed
    $\eps'>0$, and heights $h_0,\dots,h_{m-1}\ge0$. We put
    \[
      \Xi(t)=\sum_{k<m}h_k\1_{[t_k,t_{k+1})}(t),\qquad
      \omega(t)=g(t)e^{s_1t}+\Xi(t),
    \]
    and require that $\omega$ be bounded below by a positive constant on the support of $f$.
  \item A finite set $\Delta\subset[0,\infty)$ of \emph{offsets}.
\end{itemize}
We put $\vartheta_k=\frac12(t_k-\frac13)-\eps'$ (the \emph{sieve levels}), so $0<\vartheta_k<\frac12t_k$, and
\[
  I_B=\int_0^\infty\frac{e^{2st}f(t)^2}{\omega(t)}\,dt,\qquad
  D(\delta)=\int_0^\infty g(t)e^{(s_1-2\delta)t}\,dt+\sum_{k<m}h_ke^{-2\delta t_k}
    \frac{t_{k+1}^2-t_k^2}{2\vartheta_k},\qquad d_1=\frac{g(0)}6 .
\]
Here the integrand of $I_B$ is taken to be $0$ where $f=0$. We require $I_B>0$.
\end{definition}

The number $I_B$ is the first Cauchy--Schwarz factor, $D(\delta)$ bounds a diagonal
entry at offset $\delta$, and $d_1$ bounds the correlation between distinct characters.
Heath-Brown's
Lemma~12.1 is the formal analogue of the case $g=f$, $s=s_1=\lambda_1$, $\Xi=0$ and all offsets $0$ (in that
case $\omega=fe^{s_1t}$ is not bounded below on the support of $f$, and Heath-Brown writes $f=\sqrt
f\cdot\sqrt f$ instead).

\begin{definition}[Entries and responses]\label{def:entries}
An \emph{entry} is a triple $(\chi_j,\gamma_j,\delta_j)$ of a nonprincipal character modulo $q$, a
real height $\gamma_j$ and an offset $\delta_j\in\Delta$. Its \emph{response anchor} is
$s_j=s-\delta_j$, and its \emph{response} is
\[
  R_j=-\frac1\LL\sum_{n\ge1}\Lambda(n)\,\Re\Bigl(\chi_j(n)\,n^{-(1-s_j/\LL+i\gamma_j)}\Bigr)
  f\Bigl(\frac{\log n}\LL\Bigr).
\]
For a family of entries, the \emph{pair excess} of entry $j$ is
\[
  e_j=\sum_{\substack{k\ne j\\ \chi_k=\chi_j}}\Bigl(\Re G\bigl(-\sigma_{jk}+i(\gamma_j-\gamma_k)\LL\bigr)
  -d_1\Bigr)_+,\qquad \sigma_{jk}=s_1-\delta_j-\delta_k .
\]
\end{definition}

\begin{definition}[Safe anchor]\label{def:safeanchor}
For $\Gamma'>0$, the \emph{safe-anchor hypothesis} $\mathrm{SA}(s_1,\Gamma')$ says that no nonprincipal
$L(s,\chi)$ modulo $q$ has a zero $\rho$ with $|\Im\rho|\le\Gamma'$ and $\Re\rho>1-s_1/\LL$.
\end{definition}

The pair excess $e_j$ is needed only when two entries use the same character. It records
the part of their correlation that exceeds the bound $d_1$ for distinct characters.
The safe-anchor hypothesis ensures that discarded zero terms in the Gram estimate have
the required sign.

In the application $s_1\le\lambda_1$. Then $\mathrm{SA}(s_1,2l+1)$ holds for large $q$, by the
definition of $\lambda_1$ and \cref{inp:rectangle}: a zero with $|\gamma|\le2l+1\le10l$ and
$\Re\rho>1-s_1/\LL$ lies in $R(10l)$, hence in $R(l)$, so its parameter is at least
$\lambda_1\ge s_1$.

\subsection{Statement}

\begin{theorem}[Graded near-density lemma]\label{thm:graded}
Assume \cref{inp:X31,inp:X32,inp:burgess,inp:graham}, Mertens' estimate, and the divisor
bound (the number of divisors of $q$ is $O_\eps(q^\eps)$). Fix a near datum as in \cref{def:neardatum} and a
tolerance $\eta'>0$. There is a threshold $q_0$, depending only on this fixed data, with
the following property.

For $q\ge q_0$, let $J$ be a finite index set and let
$(\chi_j,\gamma_j,\delta_j)_{j\in J}$ be entries satisfying
\begin{enumerate}[(i)]
  \item $|\gamma_j|\le\Gamma$ for some $0\le\Gamma\le\LL/3$;
  \item $\mathrm{SA}(s_1,2\Gamma+1)$;
  \item the characters $\chi_j$ are pairwise distinct whenever $\Xi\not\equiv0$.
\end{enumerate}
Then, for every choice of real coefficients $a_j\ge0$,
\[
  \Bigl(\sum_ja_jR_j\Bigr)_+^2\le(1+\eta')\,I_B\Bigl[\sum_j\bigl(D(\delta_j)-d_1+e_j\bigr)a_j^2
  +(d_1+\eta')\Bigl(\sum_ja_j\Bigr)^2\Bigr].
\]
The responses $R_j$ and pair excesses $e_j$ are those of \cref{def:entries}.
The threshold $q_0$ is uniform in the number of entries, their characters and heights,
and the coefficients $a_j$.
\end{theorem}

Two special cases of the pair excess are used.
\begin{enumerate}[(i)]
  \item \emph{Separated heights.} If $\chi_k=\chi_j$ implies $|\gamma_j-\gamma_k|\LL\ge M$ for $k\ne j$,
    and $M$ is so large that $\Re G(-\sigma+iY)\le d_1$ for $|Y|\ge M$ and $\sigma$ in the finite set
    of values $\sigma_{jk}$, then $e_j=0$ (\cref{lem:decaytransform}; this holds once $g(0)>0$).
  \item \emph{A controlled pair.} If entry $j$ shares its character with exactly one other entry
    $k$, with $\delta_j=\delta_k=0$ and $|(\gamma_j-\gamma_k)\LL|$ in a range on which
    $\Re G(-s_1+iy)\le c^{\hi}$, then $e_j\le(c^{\hi}-d_1)_+$.
\end{enumerate}

\subsection{The idea of the proof}\label{subsec:nearidea}
Each response $R_j$ is a smoothed sum over primes, twisted by $\chi_j(n)n^{-i\gamma_j}$ and
evaluated at the point $1-s_j/\LL+i\gamma_j$. By the explicit formula (\cref{lem:response} below),
every zero $\rho$ of $L(s,\chi_j)$ near $1+i\gamma_j$ contributes about
$\Re F\bigl((\lambda_\rho-s_j)+i(\gamma_j\LL-\mu_\rho)\bigr)$ to $R_j$, and this is large when $\lambda_\rho$ is small.
So characters with zeros close to $1$ have large responses.

On the other hand, all the responses are correlations of one vector, the weighted primes,
with the twisted vectors $(\chi_j(n)n^{-i\gamma_j})_n$. By the Cauchy--Schwarz inequality a
combination $\sum_ja_jR_j$ is at most the length of the prime vector, which gives the factor
$I_B$, times the length of $\sum_ja_j\chi_j(n)n^{-i\gamma_j}$ measured with the weight $\omega$. The
latter is a Gram form in the characters. Its diagonal entries are the numbers $D(\delta_j)$, and
distinct characters are nearly orthogonal on the primes: their correlation is at most
$d_1=g(0)/6$. Indeed, by the explicit formula the correlation is governed by the zeros of the
quotient character $\chi_j\bchi_k$ to the right of the safe anchor, and the safe-anchor
hypothesis says that there are none. Hence only a bounded number of characters can have large
responses.

Heath-Brown's Lemma~12.1 is the case in which all entries are tested at the same anchor. The
grading gives each entry its own anchor $s_j=s-\delta_j$. An entry tested at a larger offset
$\delta_j$ has a smaller diagonal $D(\delta_j)$, because of the factor $e^{-2\delta_jt}$, so characters
whose zeros are further from $1$ are charged less. The sieve weight $\Xi$ adds Selberg-sieve
majorants of the primes to $\omega$ on the ranges $[q^{t_k},q^{t_{k+1}})$. This makes $\omega$ larger and
the first factor $I_B$ smaller. The price is a sieve contribution to the diagonal, computed with
Graham's asymptotic, and off-diagonal sums of the quotient characters over long intervals. These
are small by Burgess's estimate provided the quotient characters are nonprincipal, that is,
provided the characters are distinct.

\subsection{A Burgess bound for all nonprincipal characters}

\begin{lemma}[Burgess bound for imprimitive characters and long intervals]\label{lem:burgessall}
Assume \cref{inp:burgess}. For every $\eps>0$ there is $C$ such that for every nonprincipal
character $\chi$ modulo $q$ and all real $N\ge0$ and $H\ge1$,
\[
  \Bigl|\sum_{N<n\le N+H}\chi(n)\Bigr|\le Cq^{1/9+\eps}\min(H,q)^{2/3}.
\]
\end{lemma}

\begin{proof}
Let $\chi$ be induced by the primitive character $\chi^*$ modulo $q^*>1$, and let $r$ be the product
of the primes dividing $q$ but not $q^*$. Then $\chi(n)=\chi^*(n)\sum_{e\mid(n,r)}\mu(e)$, so the sum
is $\sum_{e\mid r}\mu(e)\chi^*(e)\sum_{N/e<m\le(N+H)/e}\chi^*(m)$. The sum only depends on the integers
in the range, so we may take the endpoints real. In each inner sum remove complete periods of
$\chi^*$, over which $\chi^*$ sums to $0$ because $\chi^*$ is nonprincipal; what remains is a sum over at
most $\min(H/e,q^*)$ consecutive integers, which by periodicity may be shifted to start at a point
$\ge1$. \Cref{inp:burgess} bounds it by $Cq^{*1/9+\eps}\min(H/e,q^*)^{2/3}\le Cq^{1/9+\eps}\min(H,q)^{2/3}$
(for an interval of length less than $1$ the bound is trivial). By the divisor bound, the number
of $e$ is $O_\eps(q^\eps)$.
\end{proof}

\subsection{Proof of \texorpdfstring{\cref{thm:graded}}{the theorem}}
Write $t_n=\log n/\LL$. For $(n,q)=1$ put
\[
  Y_n=\sum_ja_je^{-\delta_jt_n}\chi_j(n)n^{-i\gamma_j}.
\]
Since $n^{-(1-s_j/\LL+i\gamma_j)}=n^{-1}e^{st_n}e^{-\delta_jt_n}n^{-i\gamma_j}$, and $\chi_j(n)=0$
for $(n,q)>1$,
\[
  \sum_ja_jR_j=-\Re\,\frac1\LL\sum_{(n,q)=1}\frac{\Lambda(n)}n\,e^{st_n}f(t_n)\,Y_n .
\]
\emph{Cauchy--Schwarz.} Since $\omega>0$ on the support of $f$ and $\omega\ge0$ everywhere,
\begin{equation}\label{eq:CS}
  \Bigl(\sum_ja_jR_j\Bigr)_+^2\le
  \Bigl(\frac1\LL\sum_n\frac{\Lambda(n)}n\,\frac{e^{2st_n}f(t_n)^2}{\omega(t_n)}\Bigr)
  \Bigl(\frac1\LL\sum_{(n,q)=1}\frac{\Lambda(n)}n\,\omega(t_n)\,|Y_n|^2\Bigr).
\end{equation}

\emph{The first factor.} The function $\varphi_B(t)=e^{2st}f(t)^2/\omega(t)$ is bounded, because
$\omega$ is bounded below on the support of $f$; it vanishes outside the support of $f$ and is
continuous except at finitely many points, so it is Riemann integrable. Given $\eps_B>0$ choose a
partition $0=u_0<\dots<u_K$ of an interval containing the support of $f$ whose upper Riemann
sum is at most $I_B+\eps_B/2$. By Mertens' estimate $\sum_{n\le x}\Lambda(n)/n=\log x+O(1)$ (see \cite{RosserSchoenfeld1962approximate}
for explicit forms) we have
$\LL^{-1}\sum_{q^{u_i}<n\le q^{u_{i+1}}}\Lambda(n)/n=u_{i+1}-u_i+O(1/\LL)$, so the first factor is at most
$I_B+\eps_B$ for large $q$.

\emph{The second factor.} Split it as $Q_1+Q_2$ according to $\omega=ge^{s_1t}+\Xi$. Expanding
$|Y_n|^2$ and writing $\chi_{jk}=\chi_j\bchi_k$,
\[
  Q_1=\sum_{j,k}a_ja_k\,X_{jk},\qquad
  X_{jk}=\Re\frac1\LL\sum_{(n,q)=1}\Lambda(n)\chi_{jk}(n)\,n^{-(1-\sigma_{jk}/\LL)-i(\gamma_j-\gamma_k)}\,
  g(t_n).
\]
All the points $w_{jk}=1-\sigma_{jk}/\LL+i(\gamma_j-\gamma_k)$ involved satisfy $|\Re w_{jk}-1|\le C/\LL$ and
$|\Im w_{jk}|\le2\Gamma\le\LL$, so \cref{inp:X31,inp:X32} apply uniformly. The error terms
$o(1)$ below are uniform in the characters, heights and entries.
\begin{itemize}
  \item \emph{Diagonal, $j=k$.} Here $\chi_{jj}=\chi_0$ and the height is $0$. By \cref{inp:X31},
    \[
      X_{jj}=\int_0^\infty g(t)e^{(s_1-2\delta_j)t}\,dt+o(1).
    \]
  \item \emph{Distinct characters.} Here $\chi_{jk}$ is nonprincipal. By \cref{inp:X32},
    \[
      X_{jk}\le-\sum_{\rho}\Re G\bigl((w_{jk}-\rho)\LL\bigr)+\frac{\varphi(\chi_{jk})}2\,g(0)+o(1),
    \]
    the sum running over the zeros of $L(s,\chi_{jk})$ in a disc of radius $\delta<1$ about
    $1+i(\gamma_j-\gamma_k)$. Such a zero has $|\Im\rho|\le2\Gamma+1$, so by $\mathrm{SA}(s_1,2\Gamma+1)$ it
    has $\Re\rho\le1-s_1/\LL\le1-\sigma_{jk}/\LL=\Re w_{jk}$. Hence $\Re(w_{jk}-\rho)\LL\ge0$, and each term
    $\Re G(\cdot)$ is $\ge0$ by Condition~2. Dropping them and using $\varphi\le\frac13$, $g(0)\ge0$,
    we get $X_{jk}\le d_1+o(1)$. When $\sigma_{jk}<0$ the point $w_{jk}$ lies to the right of $1$, and
    the same argument applies.
  \item \emph{The same character, $j\ne k$.} Here $\chi_{jk}=\chi_0$ at the height
    $\gamma_j-\gamma_k$. By \cref{inp:X31},
    $X_{jk}=\Re G(-\sigma_{jk}+i(\gamma_j-\gamma_k)\LL)+o(1)$.
\end{itemize}
For $a\ge0$ the off-diagonal part is therefore at most
\[
  \sum_{j\ne k}a_ja_k\bigl(d_1+\epsilon_{jk}\bigr)+o(1)\Bigl(\sum_ja_j\Bigr)^2,\qquad
  \epsilon_{jk}=\1_{\chi_j=\chi_k}\bigl(\Re G(-\sigma_{jk}+i(\gamma_j-\gamma_k)\LL)-d_1\bigr)_+ .
\]
Since $\epsilon_{jk}=\epsilon_{kj}$ ($\Re G$ is even in the imaginary part, $g$ being real) and
$2a_ja_k\le a_j^2+a_k^2$, we have $\sum_{j\ne k}a_ja_k\epsilon_{jk}\le\sum_je_ja_j^2$. Also
$\sum_{j\ne k}a_ja_k=(\sum_ja_j)^2-\sum_ja_j^2$. Hence
\begin{equation}\label{eq:Q1}
  Q_1\le\sum_j\Bigl(\int_0^\infty g(t)e^{(s_1-2\delta_j)t}\,dt-d_1+e_j\Bigr)a_j^2
  +\bigl(d_1+o(1)\bigr)\Bigl(\sum_ja_j\Bigr)^2 .
\end{equation}

\emph{The sieve part $Q_2$.} It is present only when $\Xi\not\equiv0$, so the characters are distinct
and all pair excesses vanish. Fix a cell $k$, let $z_k=q^{\vartheta_k}$ and let
\[
  \xi_d=\mu(d)\frac{\log(z_k/d)}{\log z_k}\quad(d\le z_k),\qquad
  \nu_k(n)=\Bigl(\sum_{d\mid n,\ d\le z_k}\xi_d\Bigr)^2\ge0 .
\]
\begin{itemize}
  \item \emph{Majorant.} If $p$ is prime with $t_p\ge t_k$ then $p\ge q^{t_k}>z_k$, so the only
    divisor of $p$ up to $z_k$ is $1$ and $\nu_k(p)=\xi_1^2=1$. Hence
    $\Lambda(p)=(\log p)\,\nu_k(p)$. For $n$ not a prime power, $\Lambda(n)=0\le(\log n)\nu_k(n)$.
    Prime powers $p^e$ with $e\ge2$ and $t_{p^e}\ge t_k$ contribute
    $\ll\LL q^{-t_k/2}(\sum_ja_j)^2=o(1)(\sum_ja_j)^2$. So the contribution of cell $k$ to $Q_2$ is
    at most
    \[
      \frac{h_k}\LL\sum_{t_n\in[t_k,t_{k+1})}\frac{\log n}n\,\nu_k(n)\,|Y_n|^2+o(1)\Bigl(\sum_ja_j\Bigr)^2,
    \]
    where now $|Y_n|^2$ is expanded without the coprimality condition, with $\chi_j(n)=0$ for
    $(n,q)>1$.
  \item \emph{Diagonal.} The terms $j=k'$ of $|Y_n|^2$ are at most $a_j^2e^{-2\delta_jt_n}$ with
    $\chi_j\bchi_j\le\1$. By \cref{inp:graham}, $\sum_{n\le x}\nu_k(n)=x/\log z_k+O(x/\log^2z_k)$
    for $x\ge z_k$. Partial summation gives
    \[
      \frac1\LL\sum_{t_n\in[t_k,t_{k+1})}\frac{\log n}n\,\nu_k(n)e^{-2\delta_jt_n}
      =\int_{t_k}^{t_{k+1}}\frac t{\vartheta_k}e^{-2\delta_jt}\,dt+O\Bigl(\frac1\LL\Bigr)
      \le e^{-2\delta_jt_k}\frac{t_{k+1}^2-t_k^2}{2\vartheta_k}+O\Bigl(\frac1\LL\Bigr).
    \]
  \item \emph{Off-diagonal.} For $j\ne k'$ the character $\chi_{jk'}=\chi_j\bchi_{k'}$ is nonprincipal, since the
    characters are distinct.
    Put $W(n)=(\log n)n^{-1}e^{-(\delta_j+\delta_{k'})t_n}n^{-i(\gamma_j-\gamma_{k'})}$. Opening
    $\nu_k$, the term is
    \[
      \frac1\LL\sum_{d,e\le z_k}\xi_d\xi_e\,\chi_{jk'}([d,e])\sum_{m}\chi_{jk'}(m)\,W([d,e]m),
    \]
    where $m$ runs over an interval $[M_0,M_1)$ with $M_0=q^{t_k}/[d,e]\ge q^{t_k-2\vartheta_k}=
    q^{1/3+2\eps'}$. By \cref{lem:burgessall},
    $|\sum_{M_0\le m\le u}\chi_{jk'}(m)|\ll u^{2/3}q^{1/9+\eps}$ for every real $u\ge M_0$. We also have
    $|W(u[d,e])|\ll\LL(u[d,e])^{-1}$ and
    $|\partial_uW(u[d,e])|\ll\LL(1+|\gamma_j-\gamma_{k'}|)u^{-2}[d,e]^{-1}\ll\LL^2u^{-2}[d,e]^{-1}$.
    Abel summation, with the bound for the partial sums $\sum_{M_0\le m\le u}\chi_{jk'}(m)$ at every
    $u$, therefore bounds the inner sum by
    \[
      \ll\LL^2[d,e]^{-1}q^{1/9+\eps}M_0^{-1/3}=\LL^2q^{1/9+\eps-t_k/3}[d,e]^{-2/3}.
    \]
    Since $\sum_{d,e\le z_k}[d,e]^{-2/3}\le\sum_{g\le z_k}g^{-2/3}\bigl(\sum_{d'\le z_k/g}d'^{-2/3}\bigr)^2\le
    9\zeta(\frac43)z_k^{2/3}$ and $|\xi_d|\le1$, the term is
    \[
      \ll\LL\,q^{1/9+\eps-t_k/3+2\vartheta_k/3}=\LL\,q^{\eps-2\eps'/3},
    \]
    because $2\vartheta_k=t_k-\frac13-2\eps'$. With $\eps<\eps'/2$ this is $o(1)$ uniformly, and
    the sum over $j,k'$ is $o(1)(\sum_ja_j)^2$. (A bound for the partial sums only at the end
    point of the interval, combined with the total variation of $W$, would lose a factor
    $(M_1/M_0)^{2/3}$, which is too large for the cells used.)
\end{itemize}
Summing over the cells,
\begin{equation}\label{eq:Q2}
  Q_2\le\sum_ja_j^2\sum_{k<m}h_ke^{-2\delta_jt_k}\frac{t_{k+1}^2-t_k^2}{2\vartheta_k}
  +o(1)\Bigl(\sum_ja_j\Bigr)^2 .
\end{equation}

\emph{Conclusion.} By \eqref{eq:Q1} and \eqref{eq:Q2} the second factor is at most
$B+o(1)(\sum_ja_j)^2$, where
$B=\sum_j(D(\delta_j)-d_1+e_j)a_j^2+d_1(\sum_ja_j)^2\ge\sum_j(D(\delta_j)+e_j)a_j^2\ge0$. With the first
factor $I_B+o(1)$ and $I_B>0$, the right side of \eqref{eq:CS} is at most
$(1+\eta')I_B[B+\eta'(\sum_ja_j)^2]$ for $q$ large. \qed

\subsection{Responses and the threshold form}

\begin{lemma}[Uniform decay in the imaginary direction]\label{lem:decaytransform}
If $f$ satisfies Condition~1, then $\Re F(x+iY)\to0$ as $|Y|\to\infty$, uniformly for $x$ in
compact sets.
\end{lemma}

\begin{proof}
Integrating by parts twice, $F(z)=f(0)/z+O_K(|z|^{-2})$ for $\Re z$ in a compact set $K$. Moreover
$\Re(f(0)/z)=f(0)\Re z/|z|^2=O_K(|z|^{-2})$.
\end{proof}

\begin{lemma}[Response lemma]\label{lem:response}
Assume \cref{inp:X32}. Let $f$ satisfy Conditions~1 and~2, and fix $\sigma_{\max}>0$, an integer $K_0\ge0$, and
$\eta'>0$. There are $M>0$, $\delta\in(0,1)$ and $q_0$ such that the following holds for $q\ge q_0$,
and it remains true when $\delta$ is replaced by any smaller radius (with a suitable $q_0$).
Let $\chi$ be nonprincipal, $|\sigma|\le\sigma_{\max}$ and $|\gamma|\le\LL$. Let $\cS$ be a finite set
of distinct zeros of $L(s,\chi)$ in the disc $|1+i\gamma-\rho|\le\delta$, each retained with its
full multiplicity $m_\rho$. Suppose that every other zero of
$L(s,\chi)$ in the disc with $\lambda_\rho<\sigma$ satisfies $|\gamma\LL-\mu_\rho|\ge M$, and that there
are at most $K_0$ such zeros, counted with multiplicity. Then the response of $(\chi,\gamma,\sigma)$,
that is $R=-\LL^{-1}\sum_n\Lambda(n)\Re(\chi(n)n^{-(1-\sigma/\LL+i\gamma)})f(t_n)$, satisfies
\[
  R\ge\sum_{\rho\in\cS}m_\rho\,\Re F\bigl((\lambda_\rho-\sigma)+i(\gamma\LL-\mu_\rho)\bigr)
  -\frac{f(0)}6-\eta' .
\]
\end{lemma}

\begin{proof}
By \cref{lem:decaytransform} choose $M$ with $|\Re F(x+iY)|\le\eta'/(2K_0+2)$ for $|Y|\ge M$ and
$x\in[-\sigma_{\max},0]$. Apply \cref{inp:X32} with $\eps=\eta'/2$ at the point
$w=1-\sigma/\LL+i\gamma$, and let $\delta$ be its disc radius, or any smaller radius (see the remark after
\cref{inp:X32}). We have $(w-\rho)\LL=
(\lambda_\rho-\sigma)+i(\gamma\LL-\mu_\rho)$, so
\[
  R\ge\sum_{|1+i\gamma-\rho|\le\delta}\Re F\bigl((w-\rho)\LL\bigr)-\frac{\varphi(\chi)}2f(0)-\frac{\eta'}2 ,
\]
the sum counting multiplicity. Zeros in $\cS$ are kept. A zero outside $\cS$ with
$\lambda_\rho\ge\sigma$ has $\Re(w-\rho)\LL\ge0$ and a term $\ge0$ (Condition~2), which we drop. The
remaining zeros have $\lambda_\rho\in[0,\sigma)$ and $|\gamma\LL-\mu_\rho|\ge M$; there are at most $K_0$
of them, and each term is at least $-\eta'/(2K_0+2)$. Finally $\varphi(\chi)\le\frac13$ and
$f(0)\ge0$.
\end{proof}

\begin{proposition}[Threshold form]\label{prop:threshold}
Let $J$ be a finite index set, let $D_j>0$ and $v_j\in\RR$ for $j\in J$, and let $d>0$.
Suppose that
\[
  \Bigl(\sum_j a_jv_j\Bigr)_+^2\le\sum_jD_ja_j^2+d\Bigl(\sum_j a_j\Bigr)^2
  \qquad\text{for every }(a_j)_{j\in J}\in[0,\infty)^J.
\]
Then there is
$\tau\in[0,\sqrt d]$ with
\begin{equation}\label{eq:T}\tag{T}
  \sum_j\frac{(v_j-\tau)_+^2}{D_j}\le1-\frac{\tau^2}d .
\end{equation}
\end{proposition}

We call the $v_j$ the \emph{features}, the $D_j$ the \emph{diagonals} and $d$ the \emph{correlation term}
of the inequality, and $\tau$ its \emph{threshold}. A feature measures how strongly an entry is
detected, a diagonal what the entry costs on its own, and $d$ bounds the correlation between
distinct entries.

\begin{example}[One level and two levels]\label{ex:twolevels}
(i) Suppose that $N$ entries have the same feature $v>0$ and the same diagonal $D$, with $v^2>d$.
Taking every $a_j=1$ in the hypothesis of \cref{prop:threshold} gives $(Nv)^2\le ND+dN^2$, that is,
\[
  N\le\frac{D}{v^2-d}.
\]
This is the kind of bound that Heath-Brown's Lemma~12.1 provides: a bound for the number
$N(\lambda)$ of characters with a zero at distance at most $\lambda$, all counted alike.

(ii) Now take $D=1$ and $d=\frac14$, one entry with feature $v_1=1$ (a zero close to $1$) and two
entries with feature $v_2=\frac34$ (zeros further away). The count of (i), applied to each level
separately, allows this: it permits up to $D/(v_1^2-d)=\frac43$ entries with feature at least $v_1$,
and up to $D/(v_2^2-d)=\frac{16}5$ with feature at least $v_2$. But \eqref{eq:T} fails for every $\tau$:
for $0\le\tau\le\frac12$ its left side minus its right side is
\[
  (1-\tau)^2+2\bigl(\tfrac34-\tau\bigr)^2-\bigl(1-4\tau^2\bigr)=7\tau^2-5\tau+\tfrac98,
\]
whose discriminant $25-\frac{63}2$ is negative. (Equivalently, taking all three $a_j=1$ violates the
hypothesis: $(1+\frac32)^2=\frac{25}4>3+\frac94$.) So this configuration is excluded. A single
threshold couples the near entry and the farther ones, which separate counts cannot do. In
the leaf programs this coupling acts on all the bins of characters at once.
\end{example}

\begin{proof}
The function $\mathcal J(\tau)=\tau^2/d+\sum_j(v_j-\tau)_+^2/D_j$ is convex and $C^1$ on $[0,\infty)$, and
tends to $\infty$; let $\tau^*$ minimize it. If $\tau^*=0$ then $\mathcal J'(0)\ge0$, so all $v_j\le0$ and
$\mathcal J(0)=0$. Otherwise $\mathcal J'(\tau^*)=0$, that is $\tau^*=d\sum_ja_j$ with $a_j=(v_j-\tau^*)_+/D_j\ge0$.
Then $\sum_ja_jv_j=\sum_ja_j(v_j-\tau^*)+\tau^*\sum_ja_j=\sum_jD_ja_j^2+d(\sum_ja_j)^2=\mathcal J(\tau^*)$.
The hypothesis gives $\mathcal J(\tau^*)^2\le\mathcal J(\tau^*)$, so $\mathcal J(\tau^*)\le1$. Finally \eqref{eq:T} forces
$\tau^2\le d$.
\end{proof}

\begin{corollary}[Zero form]\label{cor:zeroform}
Fix a near datum whose detector $f$ also satisfies Condition~2. Fix a common bound
$K_0$ for the omitted zeros in \cref{lem:response}, and choose
$\sigma_{\max}\ge\max_{\delta\in\Delta}|s-\delta|$ with $\sigma_{\max}>0$.
Choose $\eta'>0$, $I_u\ge I_B$, and $D_u>0$, and put
\[
  \eta''=\eta'\min\{1,\sqrt{I_BD_u},D_u\}.
\]
For sufficiently large $q$, consider entries satisfying \cref{thm:graded} with tolerance
$\eta''$. For each entry $j$, suppose the retained set $\cS_j$ satisfies
\cref{lem:response} at anchor $s_j$, with the separation and disc radius chosen for
tolerance $\eta''$. Define
\[
  r_j=\sum_{\rho\in\cS_j}m_\rho\,\Re F\bigl((\lambda_\rho-s_j)+i(\gamma_j\LL-\mu_\rho)\bigr)-\frac{f(0)}6 ,
\]
and choose $D_j^+>0$ with $D_j^+\ge D(\delta_j)-d_1+e_j$ for every $j$.
Then there is $\tau\in[0,\sqrt d]$ satisfying \eqref{eq:T} with
\[
  v_j=\frac{r_j}{\sqrt{I_uD_u}}-\eta',\qquad D_j=\frac{(1+\eta')D_j^+}{D_u},\qquad
  d=\frac{(1+\eta')(d_1+\eta'D_u)}{D_u}\ \bigl(=(1+\eta')(d_1/D_u+\eta')\bigr).
\]
The same threshold remains valid if features $v_j$ are decreased, diagonals $D_j$ or
the correlation term $d$ are increased, or nonpositive features are replaced by $0$.
An entry with nonpositive feature may also be omitted.
\end{corollary}

\begin{proof}
Apply \cref{thm:graded,lem:response} with the tolerance $\eta''$ fixed in the statement.
The response lemma gives $R_j\ge k_j:=r_j-\eta''$, so for $a\ge0$
we have $(\sum_ja_jk_j)_+\le(\sum_ja_jR_j)_+$. Dividing the inequality of \cref{thm:graded} by
$I_BD_u$, we get for all $a\ge0$
\[
  \Bigl(\sum_ja_j\frac{k_j}{\sqrt{I_BD_u}}\Bigr)_+^2\le
  \sum_j\frac{(1+\eta'')D_j^+}{D_u}\,a_j^2+\frac{(1+\eta'')(d_1+\eta'')}{D_u}
  \Bigl(\sum_ja_j\Bigr)^2 ,
\]
where we also replaced each coefficient $D(\delta_j)-d_1+e_j$ by the larger $D_j^+$.
\Cref{prop:threshold} gives $\tau$ satisfying \eqref{eq:T} with these features, diagonals and
correlation term; only the positivity of the $D_j^+$ is needed, not that of $D(\delta_j)-d_1+e_j$. We compare with the displayed quantities.
\begin{itemize}
  \item If $r_j>0$ then $v_j\le k_j/\sqrt{I_BD_u}$, because $I_u\ge I_B$ and
    $\eta''/\sqrt{I_BD_u}\le\eta'$. If $r_j\le0$ then $v_j<0$.
  \item $(1+\eta')\ge(1+\eta'')$, so the displayed diagonals are at least the ones above.
  \item $(1+\eta')(d_1+\eta'D_u)\ge(1+\eta'')(d_1+\eta'')$, since $\eta''\le\eta'D_u$.
\end{itemize}
Finally, if \eqref{eq:T} holds at $\tau$ for features $v_j$, diagonals $D_j$ and correlation term $d$, it
also holds at $\tau$ for features $v'_j\le\max(v_j,0)$ with $v'_j\le v_j$ or $v'_j\le0$, diagonals
$D'_j\ge D_j$ and correlation term $d'\ge d$. Indeed each term $(v'_j-\tau)_+^2/D'_j$ is at most
$(v_j-\tau)_+^2/D_j$ (the former is $0$ when $v'_j\le0\le\tau$), $1-\tau^2/d\le1-\tau^2/d'$, and
$\tau\le\sqrt d\le\sqrt{d'}$.
\end{proof}
